% Study 1 appendix sections carried over unchanged from paper_v4 (D52: Study 1 is a descriptive case study).
\section{Case study: readings fixed in advance}
\label{app:interpretations}
% Source: PREREGISTRATION_v3 section 7, PREREGISTRATION_v3_AMENDMENT_1 section 4 (D48) and
% PREREGISTRATION_v3_AMENDMENT_2 sections 3 and 5b (D49). Keep the wording; it was planned.
Each rule was planned with the sentence we report when it applies; where a reading names a place of report (``the
abstract''), it refers to the case study's own report. Readings 1--5 are from amendment 1, readings
6--8 from amendment 2, and the rest from the original analysis plan, where they now apply to the reference ladder.
\begin{enumerate}
  \item \rung{R4}$-$\rung{R5} not positive after Holm at any registry state, for a capable model: ``For this
  model, the reviewer did not discriminate better than a lookup in the registry it was given, even when the
  registry was damaged.''
  \item \rung{R4}$-$\rung{R5} positive after Holm at a damaged state but not with the registry intact: ``The
  reviewer added value over the registry only once the registry was damaged, from [the first such state].''
  \item \rung{R4}$-$\rung{R5} negative after Holm with the registry intact: ``Adding the reviewer to a correct
  registry made review worse than the registry alone.''
  \item Deference positive after Holm: ``A wrong `in force' line removed catches this reviewer makes without
  any registry: it followed the registry over the text.'' Not significant: ``We found no evidence of
  deference''; never ``the reviewer does not defer''.
  \item Any-finding catch above right-reason catch by more than \InterpPoints{} points for a capable model at
  \rung{R0} to \rung{R2}: ``Scoring by flags alone would have credited this reviewer with [the gap] points of
  catches it made for the wrong reason.''
  \item NoRec positive after Holm and the false-alarm cost small: ``Telling the reviewer that no amendment is
  recorded, instead of that the clause is in force, recovers catches the damaged record loses.'' Not significant:
  ``We found no evidence that the wording matters''; never that it does not.
  \item NoRec positive but the false-alarm cost larger than the gain: ``The wording trades catches for false alarms''
  (both numbers reported).
  \item A primary result that keeps its sign under S1: ``It does not depend on the flawed correct answers.'' One that
  changes sign: ``It depends on them.''
  \item Fewer than \MinClauses{} distinct replaced clauses in the sealed test: no primary claim. ``With
  this many distinct clauses the test cannot support a population claim; we report per-clause results
  descriptively.''
  \item \rung{R2} at least as high as \rung{R4} for the capable models: ``The bottleneck is retrieval,
  not representation: given both versions, reviewers catch the replaced rule as well as when told its
  status.''
  \item \rung{R4}$-$\rung{R2} not significant (uncorrected, secondary since amendment 1): no ladder claim. ``Telling the reviewer the
  registry status did not reliably raise catches for these models.''
  \item The weak model at \rung{R4} within \InterpPoints{} points of its \rung{R2}: ``Below a capability
  threshold, the information does not help.'' Reported as a limit of the claim.
  \item The secondary model's \rung{R2} within \InterpPoints{} points of its \rung{R4}: ``For a stronger
  model, seeing both versions is enough; the ladder result is limited to the local panel.'' Reported in
  the abstract.
  \item ``Can't confirm'' never used by a model: ``The option was available and unused.''
\end{enumerate}
``Seeing both versions is not enough'' means that \rung{R2}$-$\rung{R0} is not significant and that
\rung{R2}'s catch is below half of \rung{R4}'s.

\section{Case study: prompts}
\label{app:prompts}
The exact prompt of every rung, rendered for one dev item by \texttt{render\_prereg\_prompts.py}, is kept with
the analysis plan (\texttt{PREREGISTRATION\_v3\_PROMPTS.md}; the plan's working file names say ``preregistration'',
but it was not publicly registered) and released with the code; we do not retype it here.
\rung{R4n} is \rung{R4}'s prompt with one change: in the registry lookup block, a status value
\texttt{in\_force} reads \texttt{no\_amendment\_recorded}. A unit test checks that the two prompts differ in that word
only, with the registry intact and damaged.

\section{Case study: per-clause results}
\label{app:perclause}
\IfFileExists{generated/tab_perclause.tex}{\begin{table*}[t]\centering\scriptsize
\caption{Right-reason catches per replaced clause (k/n faults; R5 also counts S-swap faults), capable models. The clause is the unit of every interval and test in the paper.}
\label{tab:perclause}
\resizebox{\textwidth}{!}{%
\begin{tabular}{lcccccccc}
\toprule
Replaced clause & \multicolumn{4}{c}{Qwen 14B} & \multicolumn{4}{c}{Gemma 4 26B} \\
 & R2 & R4 & R4n & R5 & R2 & R4 & R4n & R5 \\
\midrule
\texttt{DL-C3-004\#cl-3.8} & 2/4 & 4/4 & 4/4 & 8/8 & 0/4 & 4/4 & 4/4 & 8/8 \\
\texttt{DL-C3-004\#cl-9.2b} & 2/2 & 2/2 & 2/2 & 4/4 & 2/2 & 2/2 & 2/2 & 4/4 \\
\texttt{DL-C3-006\#cl-3.2} & 2/3 & 3/3 & 3/3 & 6/6 & 3/3 & 3/3 & 3/3 & 6/6 \\
\texttt{DL-C3-009\#cl-3.6} & 2/4 & 4/4 & 4/4 & 8/8 & 1/4 & 4/4 & 4/4 & 8/8 \\
\texttt{GJ-C3-001\#cl-13.5-check} & 0/4 & 4/4 & 4/4 & 8/8 & 0/4 & 2/4 & 3/4 & 8/8 \\
\texttt{GJ-C3-001\#cl-13.5-class} & 3/3 & 3/3 & 3/3 & 6/6 & 2/3 & 3/3 & 3/3 & 6/6 \\
\texttt{GJ-C3-001\#cl-5-proviso} & 0/3 & 0/3 & 0/3 & 6/6 & 0/3 & 0/3 & 0/3 & 6/6 \\
\texttt{GJ-C3-003\#cl-6.2} & 0/2 & 2/2 & 2/2 & 4/4 & 0/2 & 2/2 & 2/2 & 4/4 \\
\texttt{IN-C1-011\#cl-annex-phase2-claims} & 0/2 & 2/2 & 2/2 & 4/4 & 0/2 & 2/2 & 2/2 & 4/4 \\
\texttt{KA-C5-001\#cl-15} & 0/4 & 4/4 & 3/4 & 0/8 & 0/4 & 0/4 & 0/4 & 0/8 \\
\texttt{KA-C5-001\#cl-26} & 0/4 & 0/4 & 0/4 & 0/4 & 0/4 & 0/4 & 0/4 & 0/4 \\
\texttt{TN-C5-001\#cl-2} & 0/3 & 0/3 & 0/3 & 0/6 & 0/3 & 0/3 & 0/3 & 0/6 \\
\texttt{UP-C3-001\#cl-5.1} & 1/2 & 2/2 & 2/2 & 4/4 & 0/2 & 1/2 & 1/2 & 4/4 \\
\texttt{UP-C3-001\#cl-6.2} & 1/2 & 2/2 & 2/2 & 4/4 & 1/2 & 2/2 & 2/2 & 4/4 \\
\bottomrule
\end{tabular}}
\end{table*}
}{}
Every clause-level rate in the paper is built from these cells; the released \texttt{ladder\_per\_clause.csv} has
them for every model and rung, with controls, control alarms, can't-confirm, misnamed and invalid counts.

\section{Case study: annotation agreement}
\label{app:agreement}
On the seeded shared subset (\NAgreeItemsLinks{} links, \NAgreeItemsCases{} cases) labelled by both annotators:

\begin{center}\footnotesize\setlength{\tabcolsep}{3pt}
\begin{tabular}{llrrr}
\toprule
Sheet & Field & Agreement & $\kappa$ & PABAK \\
\midrule
Links & relation & \AgreeLinksRelation & \KappaLinksRelation & \PabakLinksRelation \\
Links & effective date & \AgreeLinksEffectiveDate & \KappaLinksEffectiveDate & \PabakLinksEffectiveDate \\
Links & wording & \AgreeLinksWording & \KappaLinksWording & \PabakLinksWording \\
Links & new clause id & \AgreeLinksNewClauseId & \KappaLinksNewClauseId & \PabakLinksNewClauseId \\
Cases & decision & \AgreeCasesDecision & \KappaCasesDecision & \PabakCasesDecision \\
Cases & blocking stage & \AgreeCasesBlockingStage & \KappaCasesBlockingStage & \PabakCasesBlockingStage \\
Cases & deciding clauses & \AgreeCasesDecidingClauses & \KappaCasesDecidingClauses & \PabakCasesDecidingClauses \\
\bottomrule
\end{tabular}
\end{center}
The new-clause-id field is low mostly because one annotator wrote sub-clauses with a dot (\texttt{cl-4.2}) where the
registry writes brackets (\texttt{cl-4(2)}); the adjudicator resolved these.

\section{Case study: deviations from the analysis plan}
\label{app:deviations}
% Source: DECISIONS.md D48, D49 and the dated notes of 2026-10-29. Add any deviation found after the run.
All of the following were recorded before any model saw a test item, except the last paragraph.
\paragraph{Amendment 1.} After the dev pilot, the damaged-registry comparisons (Q1, Q2) replaced the ladder
contrasts (\rung{R4}$-$\rung{R2}, \rung{R2}$-$\rung{R0}) as the primary family; the capable models run all
\NDamageSeedsCapable{} damage seeds; discrimination, deference and the looser scoring rules (E5) were added.
\paragraph{Amendment 2, after the human audit.}
\begin{itemize}
  \item Annotation keeps a case whose adjudicated outcome (decision, and blocking stage for a rejection) matches the key;
  cases whose deciding clauses alone differ are kept and listed. The planned rule dropped them too and would have
  kept \NStrictKept{} of \NCandidateCases{} cases.
  \item The correct-answer gate (\ControlsGatePct\% judged correct) and the scorer gate (\ScorerCheckPct\% agreement)
  were replaced by reporting: \PctControlsCorrect{} and \ScorerCheckPctObserved{} were observed. The scorer was not
  changed.
  \item Added: the \rung{R4n} arm and the NoRec contrast (primary family of \NPrimaryTestsVFive{}), sensitivity sets S1 and
  S2, the confusion-count and reference-row reporting rules, descriptive splits, and the registry audit on the random
  link sample.
  \item The annotators' meeting before adjudication was dropped; the adjudicator decided every disagreement.
  \item The analysis cutoff became 24 hours after the local run started, with blocks run in priority order.
  \item The lead author approved the amendment and stated the co-authors' agreement; not every author signed it.
\end{itemize}
\paragraph{After the run.}
Dated in the decision log. (1) The local models moved from the lead author's laptop to Kaggle GPU notebooks
(\texttt{qwen2.5:14b}) and back to the laptop for the two smaller models, for speed; each model ran on one machine,
and a 30-minute partial start on the laptop that finished no block is kept and unused. Model files were checked
against the pre-specified digests on every machine. (2) The 24-hour cutoff was fixed as a wall-clock time before any
test result was read. (3) The laptop's model server crashed once mid-request; the run raised, recorded no reply, and
resumed from its cache. (4) Additions after the plan was fixed, outside every planned test: \texttt{qwen2.5:32b} (damage seeds 0--2) and Claude
Opus~4.5 through a Kaggle-hosted proxy (seed 0) on the main arms, and a same-machine repeat of three blocks to measure
determinism. Every planned block finished for the planned models.

\section{Case study: data statement}
\label{app:datastatement}
\paragraph{Curation rationale.} Instruments on rooftop solar (net metering, capacity limits,
connection procedures, central subsidy) were collected, including superseded ones, because the study
is about replaced rules that remain available.
\paragraph{Language.} English. Six of the 111 documents also contain Hindi passages. A local open model
machine-translated these passages into English for retrieval only; the translations are poor in places, and no label
depends on them.
\paragraph{Source.} Official publications of central ministries, state electricity regulatory
commissions and distribution companies in India, retrieved from their portals. Each document is
released as a URL and a hash.
\paragraph{Cases.} Synthetic applications generated from rule functions. No personal data.
\paragraph{Annotators.} Two authors with backgrounds in computing (one an AI engineer with degrees in
artificial intelligence and computer science), neither with training in law or energy policy, and a third author
adjudicating (\S\ref{sec:annotation}).
\paragraph{Licences.} Code under Apache-2.0; cases and labels under CC BY 4.0; documents as URLs and
hashes.
